
\subsubsection{2.1 Bidirectional Human-AI Alignment}

Standard formulations of AI alignment treat human preferences as exogenously given signals that AI systems are trained to satisfy [Bai et al., 2022; Rafailov et al., 2023]. Under this view, the alignment problem is unidirectional: AI behavior is adjusted to match human specifications. A bidirectional perspective recognizes that human behavior may also be shaped by repeated interaction with AI systems, creating feedback dynamics that standard alignment frameworks do not model.

Shen et al. [2024] provide the most comprehensive review of this bidirectionality gap, covering over 400 papers from machine learning, human-computer interaction, and cognitive science. Their review identifies human behavioral adaptation to AI interaction as a research gap, noting that this direction is substantially less studied than AI-to-human effects. The review establishes the conceptual case for bidirectionality as a research area but does not provide an operational framework for empirical measurement. The present study extends this analysis by attempting an empirical test and identifying the data infrastructure constraints that such testing requires.

Vishwarupe et al. [2026] audit 16 alignment benchmarks using deployment-relevant criteria and find that none include mechanisms for measuring user-facing behavioral effects of AI exposure over time. This finding indicates that the infrastructure gap is present in widely used evaluation resources, not only in smaller or more specialized datasets.

\subsubsection{2.2 Preference Dataset Design}

Shen et al. [2024] propose a framework for evaluating preference datasets along dimensions including quality, diversity, and coverage. Their framework characterizes what properties datasets must have to support specific alignment claims. The present study identifies an additional relevant property — whether datasets include annotator-level exposure metadata — and demonstrates its practical relevance: a dataset may be large and diverse while being structurally unsuitable for measuring human-to-AI adaptation.

The Anthropic HH-RLHF dataset \citep{Baietal2022} contains hundreds of thousands of pairwise preference annotations and a split structure that on first inspection resembles an annotator exposure manipulation. The present analysis shows that this interpretation is not supported by the dataset documentation. The BeaverTails dataset \citep{Jietal2023}, with 333K+ preference pairs and dual-label structure capturing response attributes, illustrates how dataset design choices constrain the questions a dataset can answer, in this case along dimensions of harm and safety rather than annotator behavioral change.

\subsubsection{2.3 RLHF and Stylistic Feature Amplification}

RLHF optimization is known to amplify surface-level stylistic features that correlate with human approval [Stiennon et al., 2020; Rafailov et al., 2023]. As reward models are trained on pairwise comparisons and policy models are optimized against reward signals, features that reliably co-occur with preferred responses become increasingly prevalent in model outputs. Documented examples include structured enumeration, hedged assertion, and explicit step-by-step reasoning.

The AIFS construct introduced in this paper operationalizes four such feature classes — structured lists, safety prefaces, chain-of-thought markers, and hedging language — as a composite count measure. The empirical result β₁ = 0.0246 (p < 0.001) confirms that AIFS features carry preference signal in the HH-RLHF corpus, consistent with prior characterizations of RLHF-amplified stylistics.

\subsubsection{2.4 Conditional Logit Methods in Preference Modeling}

Conditional logistic regression [McFadden, 1974] is an established framework for analyzing discrete choice data. The model conditions out individual-level heterogeneity via fixed effects, making it appropriate for preference pair data where annotator-specific baseline tendencies must be controlled. The use of semantic cluster fixed effects in the present study extends this approach to multi-annotator corpora where topical context is the relevant grouping variable. At the scale of 80,342 preference pairs and 27,034 cluster fixed effects, standard quasi-Newton (BFGS) optimization was found to diverge due to Hessian ill-conditioning; this paper documents the failure and the Newton-Raphson solution.

